% Part: lambda-calculus
% Chapter: introduction
% Section: conversion

\documentclass[../../../include/open-logic-section]{subfiles}

\begin{document}

\olfileid[hi]{lam}{int}{con}
\olsection{परिवर्तन और अपचयन}

लैम्ब्डा कलन में परिवर्तन या अपचयन के तीन प्रकार हैं। जैसा हम देखेंगे,
परिवर्तन और अपचयन में अंतर इस बात का है कि ``रूपांतरण'' एकदिश है या
द्विदिश।

$\alpha$-परिवर्तन चरों के पुनर्नामकरण से संबंधित है। $\lambda x. x$ और
$\lambda y.y$ को एक ही फलन (तादात्म्य फलन) का निरूपण मानना स्वाभाविक है;
इस विचार को हम इस प्रकार सटीक रूप देते हैं:


\end{document}

